\section*{Заключение}
\addcontentsline{toc}{section}{Заключение}

В ходе выполнения пяти лабораторных работ были реализованы все
поставленные задания. В лабораторной работе~1 выполнена случайная генерация
связного ациклического ориентированного графа по распределению Вейбулла,
вычислены экцентриситеты вершин, выделены центр и диаметральные вершины,
реализован метод Шимбелла для поиска кратчайших и самых длинных маршрутов
заданной длины.
В лабораторной работе~2 реализованы топологическая сортировка по алгоритму
Фалкерсона, нахождение кратчайших путей алгоритмом Флойда-Уоршалла
с выводом матрицы расстояний и сравнение алгоритмов по числу итераций.
В лабораторной работе~3 построена сеть, найден максимальный поток
алгоритмом Форда--Фалкерсона и вычислен поток минимальной стоимости
величиной $\lfloor 2/3 \cdot F_{\max} \rfloor$.
В лабораторной работе~4 подсчитано число остовных деревьев по теореме Кирхгофа,
построен минимальный остов алгоритмом Краскала, выполнено кодирование
и декодирование кода Прюфера с сохранением весов рёбер, а также найдено
максимальное независимое множество вершин.
В лабораторной работе~5 для неориентированного графа реализована проверка того, что граф является Эйлеровым, а также достройка графа до Эйлерова На основе добавления и удаления ребер. Построен эйлеров цикл алгоритмом Хирхольцера. На основе минимального остова (алгоритм Краскала)
получена фундаментальная система циклов, и реализованна симметрическая разность которая для двух и более циклов позволяет строить новые циклы.

\subsection*{Достоинства реализации}

Все пять лабораторных работ используют один и тот же сгенерированный граф, что помогает при тестировании работы разных алгоритмов на одних
и тех же данных.

Тоже самое можно сказать про реализованные алгоритмы, которые можно переиспользовать в разных лабораторный работах. Например, алгоритм Краскала в четвертой лабораторной работе используется для построения минимального остова, а в пятой лабораторной работе для получения фундаментальной системы циклов. Это позволяет избежать дублирования кода.

Использование контейнеров STL позволяет не реализовывать базовые структуры данных вручную. Например, \texttt{vector} используется для хранения списков вершин и рёбер, \texttt{queue} и \texttt{stack} для обходов графа в BFS, \texttt{set} для хранения множеств рёбер при работе с фундаментальными циклами.
\subsection*{Недостатки реализации}

Граф хранится в виде матрицы смежности, что влияет на использованную память и время поиска вершин. Это неэффективно для графов с маленьким количеством ребер, так как все равно требуется $O(n^2)$ памяти. И при поиске соседей вершины приходится за $O(n)$ по всей строке матрицы.

В алгоритме Краскала приходится дублировать ребра в отдельную структуру данных, чтобы потом запускать DFS для проверки наличия циклов.

Нет ограничения на число вершин в графе, что может привести к ошибкам при генерации и дальнейшему падению программы.
\subsection*{Масштабирование}

Можно заменить матричное представление графа на списки смежности,
чтобы программа работала быстрее и занимала меньше памяти на разреженных
графах.

Также можно добавить ограничение на число вершин в графе, чтобы избежать ошибок при генерации и падению программы.

И можно реализовать графический интерфейс для визуализации графа и результатов алгоритмов, что сделает программу более удобной и наглядной. Например через библиотеку Qt или с использование веб-интерфейса на HTML/CSS/JS.